\section*{ЗАКЛЮЧЕНИЕ}
\addcontentsline{toc}{section}{ЗАКЛЮЧЕНИЕ}

В ходе выполнения курсовой работы была разработана программная система для моделирования продукционной системы Маркова на языке Python. Система реализована с графическим интерфейсом с использованием библиотек Tkinter и CustomTkinter, позволяет загружать правила Маркова из текстовых файлов, выполнять симуляцию преобразований строк, вести пошаговую трассировку и сохранять результаты работы в лог-файлы.

Основные преимущества разработанной системы заключаются в наглядности выполнения правил, возможности анализа статистики применения правил, удобстве конфигурации системы через файлы правил и интерактивном пользовательском интерфейсе. Система демонстрирует принципы работы продукционных систем и может использоваться в образовательных целях для изучения формальных моделей вычислений.

Основные результаты работы:

\begin{enumerate}
	\item Проведен анализ предметной области. Изучены теоретические основы продукционных систем, принципы работы правил Маркова и алгоритмы их применения.
	\item Разработана концептуальная модель системы симуляции. Определены основные сущности (правила, текущая строка, шаги выполнения), их атрибуты и связи. Разработана структура хранения правил и логов.
	\item Осуществлено проектирование системы. Разработана архитектура приложения с разделением на модули (core, gui, utils). Спроектированы классы для работы с правилами, симуляции процесса преобразований и ведения статистики.
	\item Реализована и протестирована система симуляции. Созданы классы MarkovRules, MarkovSystem и компоненты GUI. Проведено тестирование работы правил на различных сценариях, проверена корректность пошаговой симуляции и сохранения логов.
\end{enumerate}

Все требования, заявленные в техническом задании, были полностью выполнены. Разработанное приложение успешно выполняет моделирование продукционной системы Маркова, поддерживает пошаговый режим с визуализацией, отображение статистики по применению правил и ведение логов.

Готовый рабочий проект представлен в виде Python-приложения с графическим интерфейсом, готового к использованию. Приложение может быть расширено дополнительным функционалом, например, поддержкой других форматов правил или более сложной визуализации процесса преобразований.